\documentclass{article}
\usepackage[UTF8]{ctex}
\usepackage{amsmath,amssymb,amsfonts}
\usepackage{enumitem}
\usepackage{geometry}
\geometry{a4paper, margin=1in}
% === 错题本排版优化（不参与题目标识） ===
\setlength{\parindent}{0pt}
\setlength{\parskip}{0.35em plus 0.10em minus 0.08em}
\setlength{\emergencystretch}{3em}
\setlist{itemsep=0.25em, topsep=0.40em, parsep=0pt}
\allowdisplaybreaks[2]
% === 错题本排版优化结束 ===
\begin{document}
\section*{东南大学2024年数学分析真题}
\section*{86．东南大学2024年数学分析试卷}
\begin{itemize}

\subsection{题目1}
  \item[1．] 设 \(y=y(x)\) 是由方程
\[
\arcsin (x y)+y^{2}-y e^{x y}=2
\]
所确定的隐函数，求曲线 \(y=y(x)\) 在点 \((0,2)\) 处的切线方程．

\subsection{题目2}
  \item[2．] 求极限 \(\lim _{n \rightarrow \infty}\left(\sqrt[n]{n^{2}+1}-1\right) \sin \frac{n \pi}{2}\) ．
\end{itemize}

\subsection{题目3}
3．计算不定积分 \(\int \frac{\arccos x}{x^{2}} \mathrm{~d} x\) ．

\subsection{题目4}
4．计算定积分 \(\int_{-1}^{1} \frac{x+2}{e^{x}+e^{-x}} \mathrm{~d} x\) ．

\subsection{题目5}
5．求幂级数 \(\sum_{n=1}^{\infty} \frac{n+1}{n!2^{n}}(x-2)^{n}\) 的收敛域与和函数．

\subsection{题目6}
6．计算二重积分 \(\iint_{D} \frac{\sin y \cos y}{y} \mathrm{~d} x \mathrm{~d} y\) ，其中 \(D\) 为直线 \(y=x\) 与抛物线 \(x=y^{2}\) 所围成的区域。

\subsection{题目7}
7．设 \(f(x, y)=\left\{\begin{array}{ll}\frac{x^{2} y}{\sqrt{x^{2}+y^{2}}}, & x^{2}+y^{2} \neq 0 ; \\ 0, & x^{2}+y^{2}=0 .\end{array}\right.\) 求二阶混合偏导数 \(f_{x y}\) ．

\subsection{题目8}
8．设 \(f(x)=\pi-x, x \in[0, \pi]\) ．

\begin{itemize}
  \item[（1）] 将 \(f(x)\) 展开为余弦级数，并在 \([-\pi, \pi]\) 上写出和函数的表达式．
  \item[（2）] 判断该级数在 \([0, \pi]\) 内是否一致收敛，并说明原因．
\end{itemize}

\subsection{题目9}
9．计算曲线积分

\[
\int_{L} y^{2} \mathrm{~d} x+z^{2} \mathrm{~d} y+x^{2} \mathrm{~d} z
\]

其中 \(L\) 为球面 \(x^{2}+y^{2}+z^{2}=a^{2}(z \geq 0, a>0)\) 与柱面 \(x^{2}+y^{2}-a x=0\) 的交线，从 \(z\) 轴正向看去取逆时针方向。

\subsection{题目10}
10．解答如下问题：

\begin{itemize}
  \item[（1）] 求极限 \(\lim _{x \rightarrow 0^{+}} \frac{1-\cos x}{\int_{0}^{x} \frac{\ln (1+x y)}{y} \mathrm{~d} y}\) ．
  \item[（2）] 计算含参量反常积分 \(\int_{0}^{+\infty} \frac{\sin x y}{y e^{y}} \mathrm{~d} y\) ．
\end{itemize}

\subsection{题目11}
11．设 \(f(x)\) 在 \([0,1]\) 上连续，在 \((0,1)\) 内可导，

\[
f(0)=f(1)=0, f\left(\frac{1}{2}\right)=1
\]

证明：存在 \(\xi \in(0,1)\) ，使得 \(f^{\prime}(\xi)=1\) ．

\subsection{题目12}
12．设 \(D=\{(x, y) \mid 0<x<1,0<y<+\infty\}\) ，证明：对任意的 \((x, y) \in D\) ，成立不等式

\[
y x^{y}(1-x)<e^{-1}
\]

\subsection{题目13}
13．设 \(f_{n}(x)=n^{\alpha} x e^{-n x}\) ．

\begin{itemize}
  \item[（1）] \(\alpha\) 取何值时， \(\lim _{n \rightarrow \infty} \int_{0}^{1} f_{n}(x) \mathrm{d} x=\int_{0}^{1} \lim _{n \rightarrow \infty} f_{n}(x) \mathrm{d} x\) ？
  \item[（2）] \(\alpha\) 取何值时，\(\left\{f_{n}(x)\right\}\) 在 \([0,1]\) 上一致收敛？
\end{itemize}

\subsection{题目14}
14．设正项级数 \(\sum_{n=1}^{\infty} a_{n}\) 发散，以 \(S_{n}\) 表示前 \(n\) 项的和。证明：

\begin{itemize}
  \item[（1）] \(\sum_{n=1}^{\infty} \frac{a_{n}}{S_{n}}\) 发散．
  \item[（2）] \(\sum_{n=1}^{\infty} \frac{a_{n}}{S_{n}^{2}}\) 收敛。
\end{itemize}

\subsection{题目15}
15．设 \(f(x, y)\) 在区域 \(D:(0,1) \times(0,1)\) 内分别对自变量 \(x\) 和 \(y\) 连续，并且对固定的 \(x, f(x, y)\) 关于 \(y\)是单调的．证明：\(f(x, y)\) 在区域 \(D\) 内为二元连续函数．

\subsection{题目16}
16．解答如下问题：

\begin{itemize}
  \item[（1）] 叙述 \(\mathbb{R}^{n}\) 上的有限覆盖定理．
  \item[（2）] 设对任意的 \(x_{0} \in[a, b]\) ，有 \(\lim _{x \rightarrow x_{0}} f(x)=0\) ．证明：\(f(x) \in R[a, b]\) 且 \(\int_{a}^{b} f(x) \mathrm{d} x=0\) ．
\end{itemize}
\end{document}
